\documentclass[10pt,a4paper]{amsart}
\usepackage{amsmath,amssymb,amsthm,graphicx}
\usepackage{verbatim}
\newenvironment{reference}{\comment}{\endcomment}
%\newenvironment{reference}{}{}
\newenvironment{slogan}{\comment}{\endcomment}
\newenvironment{history}{\comment}{\endcomment}

% For commutative diagrams we use Xy-pic
\usepackage[all]{xy}

% We use 2cell for 2-commutative diagrams.
\xyoption{2cell}
% Theorem environments.
%
\theoremstyle{plain}
\newtheorem{theorem}[subsection]{Theorem}
\newtheorem{proposition}[subsection]{Proposition}
\newtheorem{lemma}[subsection]{Lemma}

\theoremstyle{definition}
\newtheorem{definition}[subsection]{Definition}
\newtheorem{example}[subsection]{Example}
\newtheorem{exercise}[subsection]{Exercise}
\newtheorem{situation}[subsection]{Situation}

\theoremstyle{remark}
\newtheorem{remark}[subsection]{Remark}
\newtheorem{remarks}[subsection]{Remarks}

\numberwithin{equation}{subsection}

% Macros
%
\def\lim{\mathop{\mathrm{lim}}\nolimits}
\def\colim{\mathop{\mathrm{colim}}\nolimits}
\def\Spec{\mathop{\mathrm{Spec}}}
\def\Hom{\mathop{\mathrm{Hom}}\nolimits}
\def\Ext{\mathop{\mathrm{Ext}}\nolimits}
\def\SheafHom{\mathop{\mathcal{H}\!\mathit{om}}\nolimits}
\def\SheafExt{\mathop{\mathcal{E}\!\mathit{xt}}\nolimits}
\def\Sch{\mathit{Sch}}
\def\Mor{\mathop{\mathrm{Mor}}\nolimits}
\def\Ob{\mathop{\mathrm{Ob}}\nolimits}
\def\Sh{\mathop{\mathit{Sh}}\nolimits}
\def\NL{\mathop{N\!L}\nolimits}
\def\CH{\mathop{\mathrm{CH}}\nolimits}
\def\proetale{{pro\text{-}\acute{e}tale}}
\def\etale{{\acute{e}tale}}
\def\QCoh{\mathit{QCoh}}
\def\Ker{\mathop{\mathrm{Ker}}}
\def\Im{\mathop{\mathrm{Im}}}
\def\Coker{\mathop{\mathrm{Coker}}}
\def\Coim{\mathop{\mathrm{Coim}}}

% Boxtimes
%
\DeclareMathSymbol{\boxtimes}{\mathbin}{AMSa}{"02}

%
% Macros for moduli stacks/spaces
%
\def\QCohstack{\mathcal{QC}\!\mathit{oh}}
\def\Cohstack{\mathcal{C}\!\mathit{oh}}
\def\Spacesstack{\mathcal{S}\!\mathit{paces}}
\def\Quotfunctor{\mathrm{Quot}}
\def\Hilbfunctor{\mathrm{Hilb}}
\def\Curvesstack{\mathcal{C}\!\mathit{urves}}
\def\Polarizedstack{\mathcal{P}\!\mathit{olarized}}
\def\Complexesstack{\mathcal{C}\!\mathit{omplexes}}
% \Pic is the operator that assigns to X its picard group, usage \Pic(X)
% \Picardstack_{X/B} denotes the Picard stack of X over B
% \Picardfunctor_{X/B} denotes the Picard functor of X over B
\def\Pic{\mathop{\mathrm{Pic}}\nolimits}
\def\Picardstack{\mathcal{P}\!\mathit{ic}}
\def\Picardfunctor{\mathrm{Pic}}
\def\Deformationcategory{\mathcal{D}\!\mathit{ef}}
\usepackage{cleveref}
\crefname{originallemma}{original lemma}{original lemmas}
\Crefname{originallemma}{Original lemma}{Original lemmas}
\makeatletter
\newlabel{lemma-silly}{{0F9H}{source}{}{sourcepointer-1}{}}
\newlabel{lemma-silly@cref}{{[originallemma][0][]0F9H}{[1][0][]source}{}{}{}}
\makeatother
\crefname{originaldefinition}{original definition}{original definitions}
\Crefname{originaldefinition}{Original definition}{Original definitions}
\makeatletter
\newlabel{definition-localized-chern}{{0FB5}{source}{}{sourcepointer-2}{}}
\newlabel{definition-localized-chern@cref}{{[originaldefinition][0][]0FB5}{[1][0][]source}{}{}{}}
\makeatother
\crefname{originalreference}{original reference}{original references}
\Crefname{originalreference}{Original reference}{Original references}
\makeatletter
\newlabel{situation-loc-chern}{{0GUJ}{source}{}{sourcepointer-3}{}}
\newlabel{situation-loc-chern@cref}{{[originalreference][0][]0GUJ}{[1][0][]source}{}{}{}}
\makeatother
\crefname{originallemma}{original lemma}{original lemmas}
\Crefname{originallemma}{Original lemma}{Original lemmas}
\makeatletter
\newlabel{more-algebra-lemma-lift-perfect-from-residue-field}{{0BCD}{source}{}{sourcepointer-4}{}}
\newlabel{more-algebra-lemma-lift-perfect-from-residue-field@cref}{{[originallemma][0][]0BCD}{[1][0][]source}{}{}{}}
\makeatother
\crefname{originallemma}{original lemma}{original lemmas}
\Crefname{originallemma}{Original lemma}{Original lemmas}
\makeatletter
\newlabel{lemma-loc-chern-shrink-Z}{{0FB7}{source}{}{sourcepointer-5}{}}
\newlabel{lemma-loc-chern-shrink-Z@cref}{{[originallemma][0][]0FB7}{[1][0][]source}{}{}{}}
\makeatother
\crefname{originallemma}{original lemma}{original lemmas}
\Crefname{originallemma}{Original lemma}{Original lemmas}
\makeatletter
\newlabel{lemma-loc-chern-disjoint}{{0FB6}{source}{}{sourcepointer-6}{}}
\newlabel{lemma-loc-chern-disjoint@cref}{{[originallemma][0][]0FB6}{[1][0][]source}{}{}{}}
\makeatother
\crefname{originallemma}{original lemma}{original lemmas}
\Crefname{originallemma}{Original lemma}{Original lemmas}
\makeatletter
\newlabel{lemma-loc-chern-after-pushforward}{{0FB4}{source}{}{sourcepointer-7}{}}
\newlabel{lemma-loc-chern-after-pushforward@cref}{{[originallemma][0][]0FB4}{[1][0][]source}{}{}{}}
\makeatother
\title{A localized class computed on a closed subscheme agrees with the class of a perfect object restricting to the same thing}
\author{The Stacks Project Authors}
\date{stacks-project:chow.tex:local-snapshot-20261012; GNU FDL 1.2 or later}
\begin{document}
\maketitle
\noindent\textit{Editorial note.} All delivered mathematical text is the author's own, copied from the pinned Stacks Project chapter file chow.tex; no statement or proof step was rewritten, reordered or removed.\par\medskip
\noindent\textit{Editorial note.} Preamble provenance: the pinned chapter file opens with an input of the shared preamble and carries no preamble text of its own, so the project's single Stacks-wide preamble staged at sources/stacks-chow/preamble.tex is excerpted verbatim as two spans, lines 14-18 (Xy-pic setup) and lines 154-217 (the native theorem declarations and mathematical macros). Nothing was added, removed or redefined; the wrapper class and the amsmath/amssymb/amsthm/graphicx base packages are supplied by the assembler.\par\medskip
\noindent\textit{Editorial note.} The counted statement and proof contain no cross reference, no citation and no included file, so the delivered body is self contained.\par\medskip
\noindent\textit{Editorial note.} Canonical location of the counted proposition: Stacks tag 0FB8 (sources/stacks-chow/pinned-original-source.tex line 10975).\par\medskip
\noindent\textit{Allowed standard prerequisites (editorial).}\begin{itemize}
\item The classes P'\_p and c'\_p on a closed subscheme and their characterisation. Same chapter, Lemma labelled lemma-silly, cited at proof lines 10974-10980 and 11000-11002. Supplies the characterisation that the proof verifies for the class of E.
\item The localized Chern class of a perfect object and its behaviour under shrinking the support. Same chapter, Definition labelled definition-localized-chern, invoked in the statement, and Lemma labelled lemma-loc-chern-shrink-Z, cited at proof lines 10994-10998. Writes P\_p(X\_2 -> X, E) through the class on the complement Z of U.
\item Lifting the vanishing condition from X\_1 to an open neighbourhood. More on Algebra, Lemma labelled more-algebra-lemma-lift-perfect-from-residue-field, cited at proof lines 10986-10988. Produces the open U on which E has the stated form.
\item The two defining compositions of the localized class. Same chapter, Lemma labelled lemma-loc-chern-disjoint and Lemma labelled lemma-loc-chern-after-pushforward, cited at proof lines 11004-11008. Give the vanishing against i\_\{1,*\} and the identification against i\_\{2,*\}.
\end{itemize}\medskip
\ifdefined\hypertarget\hypertarget{sourcepointer-1}{}\fi
\noindent\textit{Source reference.} Original lemma 0F9H: sources/stacks-chow/pinned-original-source.tex, source lines 9695--9695. Same-chapter pointer into the author environment lemma at source line 9695; invoked at proof line 10998 of the counted proof. It is retained only as a navigation pointer; every load-bearing step is inside the delivered span.\par\medskip
\ifdefined\hypertarget\hypertarget{sourcepointer-2}{}\fi
\noindent\textit{Source reference.} Original definition 0FB5: sources/stacks-chow/pinned-original-source.tex, source lines 10787--10787. Same-chapter pointer into the author environment definition at source line 10787; invoked at proof line None of the counted proof. It is retained only as a navigation pointer; every load-bearing step is inside the delivered span.\par\medskip
\ifdefined\hypertarget\hypertarget{sourcepointer-3}{}\fi
\noindent\textit{Source reference.} Original reference 0GUJ: sources/stacks-chow/pinned-original-source.tex, source lines 10678--10678. Same-chapter pointer into the author environment situation at source line 10678; invoked at proof line None of the counted proof. It is retained only as a navigation pointer; every load-bearing step is inside the delivered span.\par\medskip
\ifdefined\hypertarget\hypertarget{sourcepointer-4}{}\fi
\noindent\textit{Source reference.} Original lemma 0BCD: sources/stacks-chow/tags.txt, source lines 14276--14276. Cross-chapter standard prerequisite located in the project tag list staged in the local snapshot; invoked at proof line 10990 of the counted proof. It is retained only as a navigation pointer; every load-bearing step is inside the delivered span.\par\medskip
\ifdefined\hypertarget\hypertarget{sourcepointer-5}{}\fi
\noindent\textit{Source reference.} Original lemma 0FB7: sources/stacks-chow/pinned-original-source.tex, source lines 10956--10956. Same-chapter pointer into the author environment lemma at source line 10956; invoked at proof line 10995 of the counted proof. It is retained only as a navigation pointer; every load-bearing step is inside the delivered span.\par\medskip
\ifdefined\hypertarget\hypertarget{sourcepointer-6}{}\fi
\noindent\textit{Source reference.} Original lemma 0FB6: sources/stacks-chow/pinned-original-source.tex, source lines 10940--10940. Same-chapter pointer into the author environment lemma at source line 10940; invoked at proof line 11004 of the counted proof. It is retained only as a navigation pointer; every load-bearing step is inside the delivered span.\par\medskip
\ifdefined\hypertarget\hypertarget{sourcepointer-7}{}\fi
\noindent\textit{Source reference.} Original lemma 0FB4: sources/stacks-chow/pinned-original-source.tex, source lines 10918--10918. Same-chapter pointer into the author environment lemma at source line 10918; invoked at proof line 11008 of the counted proof. It is retained only as a navigation pointer; every load-bearing step is inside the delivered span.\par\medskip
\begin{lemma}
\label{lemma-loc-chern-agree}
In Lemma \ref{lemma-silly} say $E_2$ is the restriction of a perfect
$E \in D(\mathcal{O}_X)$ whose restriction to $X_1$ is zero,
resp.\ isomorphic to a finite locally free $\mathcal{O}_{X_1}$-module
of rank $< p$ sitting in cohomological degree $0$. Then the class
$P'_p(E_2)$, resp.\ $c'_p(E_2)$ of Lemma \ref{lemma-silly} agrees with
$P_p(X_2 \to X, E)$, resp.\ $c_p(X_2 \to X, E)$ of
Definition \ref{definition-localized-chern} provided $E$ satisfies
assumption (3) of Situation \ref{situation-loc-chern}.
\end{lemma}

\begin{proof}
The assumptions on $E$ imply that there is an open $U \subset X$
containing $X_1$ such that $E|_U$ is zero, resp.\ isomorphic to a finite locally
free $\mathcal{O}_U$-module of rank $< p$. See More on Algebra, Lemma
\ref{more-algebra-lemma-lift-perfect-from-residue-field}.
Let $Z \subset X$ be the complement of $U$ in $X$ endowed with
the reduced induced closed subscheme structure. Then
$P_p(X_2 \to X, E) = (Z \to X_2)_* \circ P_p(Z \to X, E)$,
resp.\ $c_p(X_2 \to X, E) = (Z \to X_2)_* \circ c_p(Z \to X, E)$
by Lemma \ref{lemma-loc-chern-shrink-Z}.
Now we can prove that $P_p(X_2 \to X, E)$, resp.\ $c_p(X_2 \to X, E)$
satisfies the characterization of $P'_p(E_2)$, resp.\ $c'_p(E_2)$
given in Lemma \ref{lemma-silly}. Namely, by the relation
$P_p(X_2 \to X, E) = (Z \to X_2)_* \circ P_p(Z \to X, E)$,
resp.\ $c_p(X_2 \to X, E) = (Z \to X_2)_* \circ c_p(Z \to X, E)$
just proven and the fact that $X_1 \cap Z = \emptyset$,
the composition $P_p(X_2 \to X, E) \circ i_{1, *}$,
resp.\ $c_p(X_2 \to X, E) \circ i_{1, *}$ is zero
by Lemma \ref{lemma-loc-chern-disjoint}.
On the other hand,
$P_p(X_2 \to X, E) \circ i_{2, *} = P_p(E_2)$,
resp.\ $c_p(X_2 \to X, E) \circ i_{2, *} = c_p(E_2)$
by Lemma \ref{lemma-loc-chern-after-pushforward}.
\end{proof}

\end{document}
